%============================================================
% Chapter 17: MEX, Distinctness & Value-Domain Problems
% STATUS: PLANNED  (generated by tools/outline.py; replace this file when the chapter is written)
%============================================================
\chapter{MEX, Distinctness \& Value-Domain Problems}
\label{ch:value_domain}

\begin{tcolorbox}[enhanced, colback=white, colframe=secondaryblue, boxrule=1pt, arc=2mm]
\textbf{\large\textcolor{primaryblue}{Questions about which values appear rather than where: MEX, the number of distinct values, last occurrences and frequency of frequencies.}}
\end{tcolorbox}

\chapterfacts{1500}{2400}{5}{10}{30}

\begin{plannedbox}[Planned Contents --- this chapter is not written yet]
\textbf{Part A --- Lessons}
\begin{enumerate}[leftmargin=*, label=\textbf{A\arabic*.}, itemsep=1pt]
    \item The MEX: computing it, maintaining it, and how it changes when elements are added or removed
    \item Distinct values in subarrays: last occurrence arrays
    \item Frequency of frequencies and the "at most $\sqrt{2n}$ distinct frequencies" fact
    \item Permutation checks: a subarray is a permutation of $1..k$
    \item Arrays over small value domains: counting arrays per value
\end{enumerate}
\textbf{Part B --- Patterns} (with candidate problems, checked when the chapter is written)
\begin{enumerate}[leftmargin=*, label=\textbf{P\arabic*.}, itemsep=1pt]
    \item \textbf{MEX}
    \item \textbf{Distinct values in subarrays}
    \item \textbf{Frequency of frequencies}
    \item \textbf{Permutation subarrays}
    \item \textbf{Small value domains}
\end{enumerate}
\end{plannedbox}
